%% kapa: exam=1 round=27 year=2016 period=1 form=A subject=real_estate subject_ko=부동산학원론 seq=7 number=87
%% source: Q-Net 감정평가사 기출문제 (cst00309) artlSeq=5169402 file=extracted/1008/5169402/제27회 감정평가사 1차 문제/제27회 감정평가사 1교시문제(A형).pdf
%% conversion: kordoc 4.24.51 -> tools/gen1.py (auto); review=proofread-sonnet; flags=none
\begin{question}{87}
부동산권리분석에 관한 설명으로 옳지 않은 것은?
\begin{choices}
  \item 권리관계를 취급하지만 재판이나 수사행위와 같이 권력행위가 아니므로 비권력적 성격을 가진다.
  \item 우리나라 등기는 관련 법률에 다른 규정이 있는 경우를 제외하고는 당사자의 신청 또는 관공서의 촉탁에 따라 행하는 신청주의 원칙을 적용한다.
  \item 부동산권리분석을 행하는 주체가 분석대상권리의 주요한 사항을 직접 확인해야 한다는 탐문주의의 원칙은 권리분석활동을 하는 데 지켜야 할 이념이다.
  \item 자료판독을 할 때 환매특약의 등기와 신탁에 관한 등기는 소유권에 관한 사항을 기록하는 부동산 등기부의 을구에서 그 기재사항을 살펴보아야 한다.
  \item 대상부동산의 권리관계를 조사ㆍ확인하기 위한 판독 내용에는 공법상 이용제한 및 거래규제의 확인ㆍ판단이 포함된다.
\end{choices}
\end{question}
\answer{4}
